% مقدمه‌ی کوتاه پیش از فصل‌ها؛ مطابق ساختار نمونه‌ی تدوین استاد.
\chapter*{مقدمه}
\addcontentsline{toc}{chapter}{مقدمه}

رشد سریع سامانه‌های هوشمند در سلامت، امکان تحلیل حجم بزرگی از داده‌های بالینی را فراهم کرده است؛ اما در عمل، داده‌ی سلامت همواره کامل، آزاد و قابل‌تجمیع نیست. پرونده‌های الکترونیک سلامت، داده‌های آزمایشگاهی و علائم حیاتی بیماران در مراکز مختلف با قالب‌های ناهمگون ثبت می‌شوند و به‌دلیل ملاحظات حریم خصوصی، دسترسی مستقیم به آن‌ها محدود است. بنابراین مسئله‌ی اصلی این پژوهش از یک نیاز عملی شکل می‌گیرد: چگونه می‌توان از داده‌ی محدود و حساس، مدلی ساخت که برای پشتیبانی تصمیم بالینی قابل استفاده و در عین حال از نظر حریم خصوصی قابل دفاع باشد؟

این پایان‌نامه بر پیش‌بینی امکان‌پذیری درمان/ترخیص در بیماران بستری در بخش مراقبت‌های ویژه تمرکز دارد. ایده‌ی محوری آن ترکیب چند رویکرد است: مدل‌های جدولی کلاسیک برای تعیین سقف عملکرد، بازنمایی توکنی علائم برای استانداردسازی داده‌های ناهمگون، یادگیری کم‌نمونه برای سناریوهای کم‌داده، و یادگیری فدرال برای شبیه‌سازی وضعیتی که داده‌ی خام میان مراکز جابه‌جا نمی‌شود. نتایج نشان می‌دهد که یادگیری کم‌نمونه خام به‌تنهایی برای داده‌ی جدولی بالینی کافی نیست، اما وقتی با بازنمایی مناسب و مدل‌های کلاسیک در یک خط لوله‌ی ترکیبی ادغام می‌شود، می‌تواند به سقف عملکرد جدولی نزدیک شود.

ساختار پایان‌نامه به‌گونه‌ای تنظیم شده است که ابتدا مسئله، انگیزه و پرسش‌های پژوهش روشن شود؛ سپس پیشینه‌ی علمی، روش‌شناسی، طراحی آزمایش‌ها و نتایج تجربی ارائه گردد. در ادامه، تحلیل داده، بحث درباره‌ی محدودیت‌ها و نتیجه‌گیری نهایی آورده می‌شود. پیوست‌ها نیز جزئیات فنی، جداول کامل، قطعه‌کدها و مستندات تکمیلی را برای بازتولید نتایج فراهم می‌کنند.

\cleardoublepage
